\section{Fundamentação Teórica}

Seção voltada a apresentar os principais conceitos necessários para entendimento do
trabalho. É imprescindível o uso correto das referências bibliográficas.
Nesta seção, o uso de figuras pode auxiliar na explicitação do conteúdo, conforme
exemplificado pela Figura 2. Caso a Figura não seja de sua autoria, indique a fonte. Use
os recursos do Word para te auxiliar no seu trabalho: (i) Inserir Legenda, (ii) Referência
Cruzada. Desta forma, você pode criar a Lista de Figuras com a função “Inserir Índice de
Ilustrações”.
Toda Figura, Tabela e Equação devem ser devidamente citadas no texto. A Figura
2, por exemplo ilustra o controle de nível em um tanque cônico. Os valores estão
apresentados na Tabela 1. Todas as Figuras, Tabelas e Equações devem ser numeradas em
ordem crescente. Para isso utilizem o recurso de LEGENGA e REFERÊNCIA
CRUZADA
(https://www.youtube.com/watch?v=AN6MbtxOuts;
https://www.youtube.com/watch?v=GV0xBt4iaXY) para fazer as citações de Figuras,
Tabelas e Equações
Veja que a Figura 2 está organizada em um Tabela com bordas não visíveis. Isto é
um “jeitinho” para facilitar a formatação.
Figura 2: Tanque com área variável.
Para tabelas e quadros, você pode utilizar os mesmos recursos para inclusão da
legenda. Desta forma, você pode criar a Lista de Tabelas com a função “Inserir Índice de
Ilustrações”. Se necessário, poder incluir uma Lista de Quadros no documento. Veja o
exemplo da Tabela 1.
3
Tabela 1: Parâmetros do tanque cônico.
Parâmetro Valor
Recomendamos que a expressão de grandezas seja realizada seguindo o Sistema
Internacional de Unidade, SI (INMETRO, 2021), utilizando o vocabulário metrológico
adequado, seguindo o VIM (INMETRO, 2012).
Para equações, utilize o “Equation Editor” do Word, conforme exemplo:
